\chapter{Cryoglobulins}

\section*{What Is This Test?}
The cryoglobulins test detects and measures immunoglobulins that precipitate from serum or plasma at temperatures below 37\textdegree{}C (body temperature) and redissolve when warmed \cite{kaushansky2021williams}. Cryoglobulins are abnormal proteins that can cause a systemic vasculitis known as cryoglobulinemia, characterized by the deposition of immune complexes in small and medium-sized vessels. Cryoglobulins are classified by the Brouet classification into three types based on their clonality \cite{brouet1974biologic}: type I consists of a single monoclonal immunoglobulin (typically IgM or IgG) and is associated with lymphoproliferative disorders such as multiple myeloma, Waldenstr\"om macroglobulinemia, and chronic lymphocytic leukemia; type II is a mixed cryoglobulin composed of a monoclonal IgM rheumatoid factor bound to polyclonal IgG, and is strongly associated with chronic hepatitis C virus (HCV) infection; type III is composed entirely of polyclonal immunoglobulins and is associated with autoimmune and infectious diseases (systemic lupus erythematosus, rheumatoid arthritis, Sj\"ogren syndrome, and various infections). The clinical syndrome of mixed cryoglobulinemia --- the ``cryoglobulinemic vasculitis'' --- classically presents with the triad of purpura, arthralgias, and weakness, and may cause glomerulonephritis and peripheral neuropathy. The test is ordered when cryoglobulinemia is suspected and to monitor disease activity.

\section*{Alternative Names}
Cryoglobulin Test; Cryocrit; Cryoglobulins, Qualitative and Quantitative; Serum Cryoglobulins; Cryoprotein Test; Cryoglobulinemic Vasculitis Panel (with complement and HCV testing)

\section*{Principle}
The detection of cryoglobulins exploits their characteristic temperature-dependent precipitation. Because cryoglobulins precipitate in the cold, the blood must be collected and kept at 37\textdegree{}C from the moment of collection, and the serum must be separated at 37\textdegree{}C before cooling. A venipuncture is performed using a pre-warmed tube, and the sample is transported to the laboratory in a thermos or other warming device at 37\textdegree{}C. The specimen is allowed to clot at 37\textdegree{}C, the serum is separated by centrifugation at 37\textdegree{}C, and the serum is then stored at 4\textdegree{}C (refrigerator temperature) for up to 7 days. If cryoglobulins are present, a white or flocculent precipitate forms. The cryocrit is the percentage of the serum volume occupied by the precipitate after centrifugation and quantifies the amount of cryoglobulin. The precipitate is then washed at 4\textdegree{}C, redissolved at 37\textdegree{}C, and analyzed by serum protein electrophoresis and immunofixation to determine the clonality and immunoglobulin class (IgG, IgM, IgA, or light chains), which assigns the Brouet type \cite{tietz2023textbook}. Complement (C4 is typically low in mixed cryoglobulinemia) and rheumatoid factor are frequently measured concurrently.

\section*{History}
Cryoglobulins were first described in 1933 by Maxwell Wintrobe and Margaret Buell, who observed an unusual protein that precipitated in the cold in a patient with multiple myeloma. The term ``cryoglobulin'' was coined by Lerner and Watson in 1947 \cite{lerner1947studies}. In 1974, Jean-Claude Brouet and colleagues at the Paris hospital Saint-Louis published the classification of cryoglobulins into three types (type I, II, and III) that remains the standard today. The association of mixed (essential) cryoglobulinemia with hepatitis C virus was established in the early 1990s, when several groups (Pascual, Ferri, Agnello) demonstrated HCV antibodies and RNA in a large proportion of patients previously labeled as having ``essential'' mixed cryoglobulinemia \cite{agnello1992role,ferri2004mixed}; this transformed the diagnosis and treatment of the disease. Treatment evolved from corticosteroids and immunosuppressants to include antiviral therapy against HCV (which can cure the cryoglobulinemia), and, in severe disease, rituximab and plasmapheresis.

\section*{Why This Test Is Ordered}
\begin{itemize}
  \item \textbf{Evaluation of suspected cryoglobulinemic vasculitis:} Palpable purpura, arthralgias, and weakness (the classic triad); or cutaneous ulcers, Raynaud phenomenon, and acrocyanosis on cold exposure
  \item \textbf{Investigation of unexplained glomerulonephritis:} Especially membranoproliferative glomerulonephritis with low C4
  \item \textbf{Evaluation of peripheral neuropathy:} Particularly a painful, distal, sensorimotor polyneuropathy
  \item \textbf{Follow-up of chronic hepatitis C infection:} Mixed cryoglobulinemia is common in chronic HCV, and the cryocrit is monitored in patients with symptoms
  \item \textbf{Evaluation of unexplained renal, skin, or joint symptoms in autoimmune disease:} (systemic lupus erythematosus, Sj\"ogren syndrome, rheumatoid arthritis)
  \item \textbf{Investigation of lymphoproliferative disorders:} When a monoclonal cryoglobulin (type I) is suspected in myeloma or Waldenstr\"om macroglobulinemia
  \item \textbf{Evaluation of cold-induced symptoms:} Livedo reticularis, digital ischemia, and cold intolerance
\end{itemize}

\section*{Is This a Primary or Secondary Test}
The cryoglobulins test is a secondary (diagnostic) test ordered when cryoglobulinemic vasculitis is suspected. It is interpreted together with serum protein electrophoresis/immunofixation, complement levels (especially C4), rheumatoid factor, and hepatitis C serology and PCR.

\section*{How This Test Is Performed}
The specimen collection is technically demanding and must be performed correctly to avoid false-negative results. Blood is drawn into a pre-warmed plain (red-top, no anticoagulant) tube or, for some laboratories, an EDTA tube, and the sample is kept at 37\textdegree{}C throughout transport (a thermos or water bath). The laboratory allows the blood to clot at 37\textdegree{}C, separates the serum at 37\textdegree{}C, and then refrigerates the serum at 2--4\textdegree{}C for 3--7 days. The presence of a precipitate indicates a positive result. The cryocrit is expressed as a percentage of the serum volume. The precipitate is characterized by immunofixation to type the cryoglobulin (type I, II, or III) and identify the immunoglobulin class. Results are typically available within 3--7 days, because the cold-precipitation step requires days of refrigeration. Several centers use turbidimetric or nephelometric quantitation after redissolution for a quantitative result.

\section*{What Preparation Is Needed}
No fasting is required. The patient must inform the healthcare provider of any known hepatitis C infection, autoimmune disease, blood dyscrasia, or monoclonal gammopathy. All medications (especially corticosteroids, immunosuppressants, rituximab, and antiviral therapy) should be disclosed, because they can reduce cryoglobulin levels. Because cryoglobulins precipitate in the cold, the phlebotomist must use a pre-warmed tube and the sample must be kept at 37\textdegree{}C --- patients cannot collect the sample at home. The test is best performed when symptoms are active, as cryoglobulin levels may be low or undetectable between flares.

\section*{Contraindications and Precautions}
There are no absolute contraindications to the blood draw. Important interpretive cautions: (1) a negative result does not exclude cryoglobulinemia if the specimen was cooled during transport, clotted, or separated below 37\textdegree{}C --- false negatives are common with improper handling; (2) cryoglobulins may cause laboratory artifacts, including falsely elevated cold agglutinin titers, pseudoleukocytosis, pseudothrombocytosis, and pseudohyperproteinemia, because the precipitating protein is counted as cells or protein by automated analyzers --- results should be re-verified on a warmed specimen; (3) transiently low levels are seen after immunosuppression, plasma exchange, or antiviral therapy; (4) cryocrit correlates only loosely with disease severity --- a patient may have symptomatic vasculitis with a low cryocrit and vice versa; (5) monoclonal cryoglobulins (type I) cause hyperviscosity syndrome when levels are high, a medical emergency; (6) low C4 is characteristic of mixed cryoglobulinemia but is not specific; (7) in hepatitis C-negative patients, a thorough search for other infectious and autoimmune triggers is required.

\section*{Risks}
Risks are those of routine venipuncture: mild pain, bruising, small hematoma at the collection site, lightheadedness or vasovagal syncope, and extremely rarely, infection or nerve injury. No additional risk is associated with the test itself.

\section*{What Happens After the Test}
If the cryoglobulins test is positive, the cryocrit and type guide management. Type I cryoglobulinemia (monoclonal) prompts evaluation for an underlying lymphoproliferative disorder (serum protein electrophoresis, bone marrow biopsy) and treatment of the underlying disease; severe hyperviscosity is treated emergently with plasmapheresis. Mixed cryoglobulinemia (types II and III) prompts evaluation for hepatitis C (HCV antibody and RNA), autoimmune disease, and other infections. In HCV-associated cryoglobulinemia, antiviral therapy (direct-acting antivirals) can eradicate the virus and resolve the cryoglobulinemia \cite{gragnani2016prospective}. For symptomatic vasculitis, treatment includes corticosteroids, rituximab, and, in severe or rapidly progressive disease (cryoglobulinemic glomerulonephritis), plasmapheresis \cite{terrier2012management}. Complement (C4), rheumatoid factor, kidney function, and urinalysis are followed serially. The patient is counseled to avoid cold exposure, which can precipitate crises.

\section*{Understanding Results}

\subsection*{Positive (Cryoglobulins Detected)}
\begin{itemize}
  \item \textbf{Type I (monoclonal):} Associated with multiple myeloma, Waldenstr\"om macroglobulinemia, chronic lymphocytic leukemia, and monoclonal gammopathy of undetermined significance; causes hyperviscosity, cold-induced Raynaud, purpura, and digital ischemia
  \item \textbf{Type II (mixed, monoclonal IgM rheumatoid factor + polyclonal IgG):} Strongly associated with chronic hepatitis C; causes purpura, arthralgias, weakness, membranoproliferative glomerulonephritis, and peripheral neuropathy
  \item \textbf{Type III (mixed, polyclonal):} Associated with systemic lupus erythematosus, rheumatoid arthritis, Sj\"ogren syndrome, and infections (hepatitis B, HIV, Epstein-Barr virus, cytomegalovirus); generally causes milder disease
  \item \textbf{Low-level positive cryoglobulins without symptoms:} May be detected incidentally in chronic HCV or autoimmune disease and does not by itself require treatment
\end{itemize}

\subsection*{Negative}
\begin{itemize}
  \item \textbf{No cryoglobulinemia:} A negative result on a properly handled specimen makes cryoglobulinemia unlikely
  \item \textbf{False-negative due to improper handling:} Specimen cooled during transport or separated at room temperature will fail to precipitate; repeat testing with proper 37\textdegree{}C handling is warranted if suspicion remains high
  \item \textbf{Treated or quiescent disease:} Levels may be undetectable during remission or after immunosuppressive, antiviral, or plasma-exchange therapy
\end{itemize}

\section*{Normal Results}
\begin{itemize}
  \item \textbf{Negative:} No cryoglobulin precipitate is detected in serum after 7 days at 4\textdegree{}C
  \item \textbf{Cryocrit:} Normally 0--1\%; a cryocrit $>$1\% is generally considered abnormal, though low-level positive results may be clinically insignificant
  \item Reference ranges vary by laboratory; the qualitative result (present/absent) and the immunofixation type are more clinically informative than the quantitative value
\end{itemize}

\section*{Follow-up Tests}
\begin{itemize}
  \item \textbf{Serum protein electrophoresis and immunofixation:} To type the cryoglobulin and detect an associated monoclonal gammopathy
  \item \textbf{Hepatitis C antibody and HCV RNA:} Mandatory in all patients with mixed cryoglobulinemia
  \item \textbf{Hepatitis B serology and HIV testing:} For other infectious associations
  \item \textbf{Complement levels (C3, C4, CH50):} Low C4 is characteristic of mixed cryoglobulinemia
  \item \textbf{Rheumatoid factor:} Often markedly elevated as the IgM component of mixed cryoglobulins
  \item \textbf{Antinuclear antibody (ANA) and anti-dsDNA:} To screen for systemic lupus erythematosus
  \item \textbf{Urinalysis with microscopy and serum creatinine:} To detect glomerulonephritis
  \item \textbf{Quantitative immunoglobulins:} For immunoglobulin class quantitation
  \item \textbf{Bone marrow biopsy:} When type I cryoglobulinemia or an underlying lymphoproliferative disorder is suspected
  \item \textbf{Renal biopsy:} When glomerulonephritis is present and diagnosis requires confirmation
\end{itemize}

\section*{References}
\bibliographystyle{unsrtnat}
\bibliography{references}

\clearpage
